% Ad Familiares 8.17 (ad-familiares-08-17) --- English translation
% Translated by Claude (Anthropic) from the Latin (Perseus).
% Style guide: STYLE.md.

\ciceroLetterOpener{CAELIVS CICERONI S.}

\ciceroSection{1} So I would rather have been in Spain back then than at Formiae, when you set off to join Pompey! If only [corrupt] Appius Claudius were on that side, or C. Curio --- the man whose friendship dragged me, step by step, into this ruined business. For I feel that anger and affection between them have carried off my good judgement. And you: you, again, when on my way out I came to you at Ariminum by night, while you were charging me with errands of peace to Caesar and playing the marvellous citizen --- you neglected the duty of a friend, and you did not look out for me.

\ciceroSection{2} And I do not say this because I have no faith in this cause: but, believe me, it is better to die than to look at \emph{those} people. If there were not the fear of your side's cruelty, we should have been thrown out of here long ago: for on the spot, apart from a handful of money-lenders, there is not one man and not one order that is not Pompeian. For my own part, by now, I have brought it to this: that the common people, especially, and the populace that was once ours, are now yours. ``Why is this?'' you ask. Wait for the rest: I shall make you win whether you like it or not. [corrupt: ``call me a second Cato of Arruntanus''] --- you are asleep, and you don't yet seem to me to grasp where we are exposed and where we are weak. And I shall do this in hope of no reward, but for what has always counted with me most, the sake of pain and indignation. What are you doing out there? Are you waiting for a battle, the one thing you are strongest in? [corrupt] I do not know your forces; ours have got into the habit of fighting hard, of bearing cold easily, and of going hungry.
